% Workshop only: the long paper already defines these metrics in its main text.
\section{Formal Metric Definitions}
\label{app:metrics}
\label{sec:collapse-metrics}

For prompt $x$ with specification $s_x$, the validator maps an output to a
canonical valid key or failure:
$V(\cdot,s_x):\mathcal Y\to\mathcal C_x\cup\{\bot\}$.
The learner sees keys only for its generated responses, never the full
catalogue $\mathcal C_x$ or its size. For $Y\sim\pi_\theta(\cdot\mid x)$, define
\[
 P_\theta^+(x)=\Pr[V(Y,s_x)\ne\bot],\qquad
 p_\theta^+(c\mid x)=\Pr[V(Y,s_x)=c\mid V(Y,s_x)\ne\bot].
\]
The conditional distribution is defined only when $P_\theta^+(x)>0$;
otherwise all sampling metrics below are zero. Mode collapse means
concentration onto fewer valid keys, which can coexist with increasing
success probability.

For $K$ independent responses, let
\[
 D_K=\bigl|\{V(Y_i,s_x):i\le K\}\setminus\{\bot\}\bigr|,
 \qquad F_K=K^{-1}\sum_i\1[V(Y_i,s_x)\ne\bot].
\]
Then \texttt{distinct@}$K=\E[D_K]$ counts expected unique verified keys in
$[0,K]$, whereas \texttt{pass@}$K=\Pr[D_K>0]$ and
\texttt{mean@}$K=\E[F_K]$ are probabilities in $[0,1]$.
The samplewise inequality $F_K\le\1[D_K>0]\le D_K$ gives
$\texttt{mean@}K\le\texttt{pass@}K\le\texttt{distinct@}K$,
but changes in the count and probability have different units.
For a fixed prompt, $\mu_c(x)=\Pr[V(Y,s_x)=c]$ gives
\[
 \texttt{pass@}K(x)=1-(1-P_\theta^+(x))^K,\qquad
 \texttt{distinct@}K(x)=\sum_{c\in\mathcal C_x}[1-(1-\mu_c(x))^K].
\]
These identities apply before averaging over prompts. Raw
\texttt{distinct@}$K$ measures sampled coverage, not total
positive-probability support. Uniform independent sampling from $N$ finite
actions, with $n_c$ aliases of valid key $c$, gives
$\sum_c[1-(1-n_c/N)^K]$ expected keys; open-ended program spaces have no
canonical uniform generator.

The terminal co-primary outcomes use $K=8$: success probability
\texttt{pass@8} and raw \texttt{distinct@8}. Their difference
$\E[D_K-\1[D_K>0]]$ counts extra verified modes beyond the first. This
interpretive decomposition remains coupled to correctness and is neither
a third primary endpoint nor a success-matched diversity measure.
